% My choice: :
%-----------------
%\documentclass[a4paper,handout]{beamer}
%\documentclass[a4paper,presentation]{beamer}
\documentclass[aspectratio=169]{beamer}
%----------------------
%
%
%
%
%
%
%----------------------
% Yves Bouteloup :
%----------------------
%\documentclass[8pt,a4paper,presentation]{beamer}
%--------------------
% ENS templates:
%--------------------
%\documentclass{beamer}
%----------------------
%\documentclass[10pt,a4paper,presentation]{beamer}
%\documentclass[t,handout]{beamer}
%----------------------


%------------------------------
% pour windows (Font ecso0800 not found) :
%------------------------------
\usepackage{lmodern}

%-----------------------------
% The Frankfurt theme: (with )
%-----------------------------
\usetheme{Frankfurt}

%-----------------------
% Usefull packages:
%-----------------------
\usepackage{pgf}
\usepackage{bm}
%%\usepackage[french]{babel}
%%\usepackage[english,frenchb]{babel}
\usepackage[french,english]{babel} % prendre celui là ?

% Pour avoir la date et les titres en français :
%\usepackage[francais,english]{babel}
%\usepackage[latin1]{inputenc}
%\usepackage[T1]{fontenc}

%------------------------------
% pour avoir les equations :
%------------------------------
%\usepackage{sansmathaccent}
%\pdfmapfile{+sansmathaccent.map}

%----------------------------
% from ENS templates:
%----------------------------
\usepackage{pgfarrows,pgfnodes,pgfautomata,pgfheaps,pgfshade}
\usepackage{amsmath,amssymb}
%\usepackage[latin1]{inputenc}
\usepackage{colortbl} % ???
\usepackage{lipsum}
%----------------------------
% from Karim Yessad:
%----------------------------
%\usepackage[latin1]{inputenc}
%\usepackage[cyr]{aeguill}
%\usepackage[francais]{babel}

%----------------------------
% from Yves bouteloup:
%----------------------------
\usepackage[utf8]{inputenc} % !!! Important pour les accents français (Météo) !!!
\usepackage[T1]{fontenc} % !!! Important pour les accents français (Météo) !!!
%\usepackage{times} % pas certain...
\usepackage{amsfonts,amsbsy}
%\usepackage{euscript} % ??
%\usepackage{ulem} % ??
\usepackage{caption} % ??
%
\usepackage{tikz}
\usepackage{color,graphicx,subfigure}
\usepackage{epstopdf}
\usepackage{array}
\usepackage{verbatim}
%\usepackage[all]{xy}  % x-y pic (pour faire des dessins !)
\setbeamertemplate{sections/subsections in toc}[circle]
% Pour pouvoir superposer des choses (en particulier des légendes ...
\usepackage[absolute,showboxes,overlay]{textpos}  % déclaration du package
\textblockorigin{0pt}{0pt} % origine des positions (coin gauche) attention, axe y vers le bas
%\TPshowboxestrue          % affiche le contour des textblock
\TPshowboxesfalse          % n'affiche pas le contour des textblock

\DeclareMathAlphabet{\mathpzc}{OT1}{pzc}{m}{it}
%\newcommand{\lp}{\left(}
%\newcommand{\rp}{\right)}
%\newcommand{\lc}{\left[}
%\newcommand{\rc}{\right]}
%\newcommand{\la}{\left \{ }
%\newcommand{\ra}{\right \} }
\newcommand{\dr}{\partial}
\newcommand{\gray}[1]{\textcolor{gray}{#1}}
\newcommand{\red}[1]{\textcolor{red}{#1}}
\newcommand{\blue}[1]{\textcolor{blue}{#1}}
\newcommand{\mynode}[4]{\draw(#1,#2) node[right] (#3) {\phantom{AAAA}} node[right] {#4}}
% http://latexcolor.com/
\definecolor{mydarkred}{rgb}{0.5, 0.0, 0.0}
\definecolor{mydarkgreen}{rgb}{0.0, 0.5, 0.0}
\definecolor{mydarkblue}{rgb}{0.0, 0.0, 0.5}
\newcommand{\green}[1]{\textcolor{mydarkgreen}{#1}}

%-------------------
% from mf2016:
%-------------------
\definecolor{airmassmf1}{rgb}{1.000,0.941,0.686}
\definecolor{airmassmf2}{rgb}{0.768,0.756,0.690}
\definecolor{airmassmf3}{rgb}{0.568,0.607,0.682}
\definecolor{airmassmf4}{rgb}{0.384,0.482,0.655}
\definecolor{airmassmf5}{rgb}{0.227,0.388,0.623}
\definecolor{airmassmf6}{rgb}{0.086,0.313,0.588}
\definecolor{mftabname1}{rgb}{0.117,0.294,0.603}
\definecolor{mftabname2}{rgb}{0.941,0.208,0.184}
\definecolor{mftabedge1}{rgb}{0.768,0.800,0.906}
\definecolor{blockbg1}{rgb}{0.900,0.9,0.9}
\definecolor{mfblue}{rgb}{0.000,0.5607,0.7411}
%-------------------
% from mf2016:
%-------------------
%\setbeamercolor*{section in head/foot}{parent=palette tertiary}
%\setbeamercolor*{subsection in head/foot}{parent=palette primary}
%-------------------
% from mf2016:
%-------------------
% pour les pieds de page :
%\setbeamerfont{author in head/foot}{family=\sffamily}
%\setbeamerfont{title in head/foot}{family=\sffamily,series=\bfseries,size={\fontsize{6}{7}}}
%\setbeamerfont{date in head/foot}{family=\sffamily}
%\setbeamerfont{frame in head/foot}{family=\sffamily,series=\bfseries,size={\fontsize{6}{7}}}
%\setbeamerfont{section in head/foot}{size=\tiny,series=\normalfont}
%\setbeamerfont{subsection in head/foot}{family=\sffamily,size={\fontsize{1}{1}}}
%-------------------
% from mf2016: (style de couleur des item )
%-------------------
\setbeamercolor{itemize item}{fg=red} 
% ------------------------------------------------------------------------
% from mf2016: (Couleurs speciales)
% ------------------------------------------------------------------------
\definecolor{beamer@blendedblue}{rgb}{0.000,0.439,0.753}
\definecolor{mfgrey}{rgb}{0.250,0.250,0.250}
\definecolor{mfwhite}{rgb}{1.000,1.000,1.000}
\definecolor{mfblack}{rgb}{0.000,0.000,0.000}
\definecolor{mftitleblue}{rgb}{0.000,0.5607,0.7411}
\definecolor{mflogoblue}{rgb}{0.1215,0.3333,0.5843}
% Couleurs
%\setbeamercolor{frametitle}{fg=beamer@blendedblue}
%\setbeamercolor{alerted text}{fg=red}
%\setbeamercolor{block title}{use=structure,fg=white,bg=structure.fg!80!black}
%\setbeamercolor{block title alerted}{use=alerted text,fg=white,bg=alerted text.fg!80!black}
%\setbeamercolor{block body}{bg=mftabedge1}
\setbeamercolor{block body}{bg=blockbg1}
\setbeamercolor{block title example}{use=example text,fg=white,bg=example text.fg!80!black}
\setbeamercolor{block title}{fg=white,bg=mflogoblue}
\setbeamercolor{block title alerted}{use=alerted text,fg=white,bg=alerted text.fg!80!black}
\setbeamercolor{whitetext}{fg=white,bg=beamer@blendedblue!80!black}
\setbeamercolor{alerted text}{fg=red}
\setbeamercolor{frametitle}{fg=mfblue!70!mfblue, bg=white}
%\setbeamercolor{frametitle}{fg=mftitleblue}
\setbeamercolor{framesubtitle}{fg=mftitleblue}
%\setbeamercolor{block title}{use=structure,fg=white,bg=structure.fg!80!black}
%\setbeamercolor{block title alerted}{use=alerted text,fg=white,bg=alerted text.fg!80!black}
%\setbeamercolor{block title example}{use=example text,fg=white,bg=example text.fg!80!black}
%\setbeamercolor{whitetext}{fg=white,bg=beamer@blendedblue!80!black}

% ------------------------------------------------------------------------
% from mf2016: (Couleurs page de titre)
% ------------------------------------------------------------------------
%\setbeamercolor{title}{fg=mfblack}
%\setbeamercolor{subtitle}{fg=mfblack}
%\setbeamercolor{author}{fg=mftitleblue}
%\setbeamercolor{institute}{fg=mftitleblue}
%\setbeamercolor{date}{fg=mftitleblue}
% ------------------------------------------------------------------------
% from mf2016: (Polices d'ecriture)
% ------------------------------------------------------------------------
\setbeamerfont{frametitle}{family=\sffamily,series=\bfseries,size={\fontsize{13}{15}}}
\setbeamerfont{framesubtitle}{size=\normalsize}
\setbeamerfont{subsection in toc}{size=\small}
\setbeamerfont{subsubsection in toc}{size=\small}
% Polices d'ecriture page de titre
\setbeamerfont{title}{     family=\sffamily,series=\bfseries,size={\fontsize{13}{15}},parent=structure}
\setbeamerfont{subtitle}{  family=\sffamily,series=\bfseries,size=\normalsize,parent=title}
\setbeamerfont{author}{    family=\sffamily,series=\bfseries,size={\fontsize{10}{12}}}
\setbeamerfont{institute}{ family=\sffamily,series=\bfseries,size={\fontsize{10}{12}}}
\setbeamerfont{date}{      family=\sffamily,series=\bfseries,size={\fontsize{10}{12}}}

% ------------------------------------------------------------------------
% from mf2016:  (Blocs arrondis)
% ------------------------------------------------------------------------
%  \setbeamertemplate{blocks}[rounded][shadow=true,bg=yellow]
\newcommand{\bb}[1]{\begin{beamerboxesrounded}[shadow=true]{#1}}
\newcommand{\eb}{\end{beamerboxesrounded}}

% -----------------------
% from Y. Bouteloup:
% -----------------------
%\setbeamertemplate{sections/subsections in toc}[circle]
%
% Pour pouvoir superposer des choses (en particulier des légendes ...
%\usepackage[absolute,showboxes,overlay]{textpos}     % déclaration du package
%\textblockorigin{0pt}{0pt}                           %origine des positions (coin gauche) attention, axe y vers le bas !!
%\TPshowboxestrue                                  % affiche le contour des textblock
%\TPshowboxesfalse                                % n'affiche pas le contour des textblock


% ------------------------------------------------------------------------
%  from mf2016: Enlever la barre de navigation
% ------------------------------------------------------------------------
\setbeamertemplate{navigation symbols}{} 

% ------------------------------------------------------------------------
%  from mf2016: Affichage du titre
% ------------------------------------------------------------------------
%\setbeamertemplate{frametitle}
%  { %
%  \vskip0.1cm
%  \begin{beamercolorbox}[wd=\paperwidth,ht=1.1cm,leftskip=1.3cm]{frametitle}
%    \usebeamerfont{frametitle}\usebeamercolor[fg]{frametitle}\insertframetitle
%    %\vskip0.5em
%    \flushright \usebeamerfont{framesubtitle}\usebeamercolor[fg]{framesubtitle}{\itshape %\insertframesubtitle\hskip0.4cm}
%  \end{beamercolorbox}
%  \vskip0.1cm
%}
% \centering \usebeamerfont{framesubtitle}\usebeamercolor[fg]{framesubtitle}{\itshape \insertframesubtitle\hskip0.4cm}
% \flushright \usebeamerfont{framesubtitle}\usebeamercolor[fg]{framesubtitle}{\itshape \insertframesubtitle\hskip0.4cm}

% ------------------------------------------------------------------------
% from mf2016:  Reglage des listes et enumerations
% ------------------------------------------------------------------------
\setbeamertemplate{sections/subsections in toc}[square]
\setbeamertemplate{items}[default]
\setbeamertemplate{itemize item}[square]
\setbeamertemplate{itemize subitem}[circle]
\setbeamertemplate{itemize subsubitem}[triangle]
\setbeamercolor{itemize item}{fg=mftitleblue}
\setbeamercolor{itemize subitem}{fg=mfblack}
\setbeamercolor{itemize subsubitem}{fg=mfblack}

%-------------------
% from test04:       Navigation bar
%-------------------
\makeatletter
\def\insertnavigation#1{%
\vbox{{%
\usebeamerfont{section in head/foot}\usebeamercolor[fg]{section in head/foot}%
\beamer@xpos=0\relax%
\beamer@ypos=1\relax%
\hbox to #1{\hskip.3cm\setbox\beamer@sectionbox=\hbox{\kern1sp}%
  \ht\beamer@sectionbox=1.875ex%
  \dp\beamer@sectionbox=0.75ex%
    \hskip.3cm%
    \global\beamer@section@min@dim\z@
    \dohead%
    \beamer@section@set@min@width
  \box\beamer@sectionbox\hfill\hskip.3cm}%
}}}

%-------------------
% from test04:       Navigation bar / headline
%-------------------
%
\setbeamercolor{upper separation line head}{bg=mfblue!70!mfblue}
%
\setbeamertemplate{headline}
{%
\pgfuseshading{beamer@barshade}%
\ifbeamer@sb@subsection%
\vskip-9.75ex%
\else%
\vskip-6ex%
\fi%
\begin{beamercolorbox}[ignorebg,ht=2.25ex,dp=3.75ex]{section in head/foot}
\insertnavigation{\paperwidth}
\end{beamercolorbox}%
\ifbeamer@sb@subsection%
\begin{beamercolorbox}[ignorebg,ht=2.125ex,dp=1.125ex,%
  leftskip=.3cm,rightskip=.3cm plus1fil]{subsection in head/foot}
  \usebeamerfont{subsection in head/foot}\insertsubsectionhead
\end{beamercolorbox}%
\fi%
\begin{beamercolorbox}[colsep=1.5pt,ht=.75ex]{upper separation line head}
\end{beamercolorbox}
}%
\makeatother


%-----------------------------------
% from Ecrad (Y. Bouteloup):
%-----------------------------------
%\pgfdeclareimage[height=90mm,width=120mm]{garde}{cover_mf2016+cnrs+leffe}
%\pgfdeclareimage[height=96mm,width=128mm]{template}{diapo_type_mf2016+cnrs+leffe}

%-------------------
% from test04:
%-------------------
\setbeamercolor{section in head/foot}{fg=mflogoblue, bg=white}
%\setbeamercolor{section in head/foot}{fg=black, bg=white}
%
%\setbeamercolor{frametitle}{fg=red!70!black, bg=white}
%\setbeamercolor{upper separation line head}{bg=red!70!black}
%\setbeamercolor{section in head/foot}{fg=mfblue, bg=white}

\makeatletter
%------------------------------------------------------------------------------------
% from mf2016:  Pied de page (avec page et short-title, mais sans date)
%------------------------------------------------------------------------------------
%\defbeamertemplate*{footline}{Frankfurt theme}

%\setbeamertemplate{footline}
%{
%  \leavevmode%
%  \hbox{%
%  \begin{beamercolorbox}[wd=.20\paperwidth,ht=.3cm,dp=.3cm,left,leftskip=0.45cm]{frame in head/foot}%
%  %===== HERE ====================================================
%%    \usebeamerfont{section in head/foot} \usebeamercolor[mflogoblue]{section in head/foot} Page \insertframenumber / \inserttotalframenumber
%    \usebeamerfont{section in head/foot} \usebeamercolor[mflogoblue]{section in head/foot} Page \insertframenumber
%  %===== HERE ====================================================
%  \end{beamercolorbox}%
%  \begin{beamercolorbox}[wd=.60\paperwidth,ht=.3cm,dp=.15cm,center]{frame in head/foot}%
%  %===== HERE ====================================================
%    \usebeamerfont{section in head/foot} \usebeamercolor[mflogoblue]{section in head/foot} 
%    2-column model of conv. drafts accounting for the pert. pressure term
%  %===== HERE =====================================================
%  \end{beamercolorbox}%
%  }%
%  \vskip0pt%
%}


%===========================================

%-----------------------------------
% Pour les titre / sous-titre / date / afiliations / ... 
%-----------------------------------
%===== HERE ====== HERE ========
\title{
{\bf DDH: Diagnostics in Horizontal Domains} \\
\vspace{2mm}
{\bf The ddhtoolbox : a comprehensive physics-dynamics } \\
\vspace{2mm}
{\bf budget tool in ARPEGE, AROME, HARMONIE }
}
%===== HERE ====== HERE ========
%\subtitle{{\it Réunion PROC/Phys (5 décembre 2017)}}
\author{
Jean-Marcel Piriou, Fabrice Voitus, Marie Mazoyer \\
\vspace{4mm}
{\it Météo-France / CNRM / GMAP} \\
\vspace{8mm}
{\it Workshop AEMET, Madrid, 21 November 2024}
}
%===== HERE ====== HERE ========
%\institute{M\'et\'eo-France et CNRS-UMR3589 / Toulouse / France}
%\date{14 octobre 2017}
%\date{\today}
%\date{Réunion PROC/Phys (11 décembre 2017)}
\date{ }
%===== HERE ====== HERE ========

% \author[Hemaspaandra, Mukherji, Tantau]{%
%   Lane~Hemaspaandra\inst{1} \and Proshanto~Mukherji\inst{1} \and  Till~Tantau\inst{2}}
%  \institute[Universities of Rochester and Berlin]{
%    \inst{1}%
%    Department of Computer Science\\
%   University of Rochester
%    \and
%    \inst{2}%
%    Fakultät für Elektrotechnik und Informatik\\
%   Technical University of Berlin}

%============================================================================
\begin{document}


%============================================================================
\newcommand{\logos}[0]{
\begin{textblock*}{3cm}(13.4cm,8.0cm) 
\includegraphics[width=0.6cm]{logo_couleur_cnrm3_LR.eps}
\hspace{0mm} 
\includegraphics[width=0.6cm]{LOGO_MTO1_bleu_new_BR.eps}
\hspace{0mm} 
\includegraphics[width=0.6cm]{logoCNRS_UMR3589_BR.eps}  
\end{textblock*}
}
%============================================================================

%----------------------------------------------
% Ajustement du comptage des cadres
%----------------------------------------------
\addtocounter{framenumber}{-1} % La page de garde est la page zéro

%===============================================================

%=============
 \begin{frame}[plain]
%=============
  \titlepage
%=============
\begin{textblock*}{6cm}(11.0cm,7.3cm)
\includegraphics[width=1.2cm]{logo_couleur_cnrm3.eps}
\hspace{2mm}
\includegraphics[width=1.2cm]{LOGO_MTO1_bleu_new.eps}
\hspace{2mm}
\includegraphics[width=1.2cm]{logoCNRS_UMR3589.eps}
\end{textblock*}
\end{frame}

\include{com} % commandes mathématiques raccourcies, exemple: \DerP pour une dérivée partielle.


%%%%%%%%%%%%%%%%%%%%%%%%%%%%%%%%%%%%%%%%%%%%%%%%%%%%%%%%%%%%%%%%%%%%%%%%%%%%%
\section{Introduction}

%============================================================================
\subsection{Sommaire}
%============================================================================
\begin{frame}
%\frametitle{\vspace{0.5cm} \centerline{{\LARGE Sommaire}} }
\frametitle{\vspace{0.2cm} \hspace{3.5cm} \centerline{\LARGE Summary } }

\begin{textblock*}{13.5cm}(1.5cm,2.5cm)
\begin{columns}
\column{13.5cm}
\begin{block}<-1>{}
\begin{itemize}
% \Huge \huge \ LARGE \Large \large
% \normalsize \small \footnotesize \scriptsize \tiny
\item Motivation: why use DDH ? Examples of use.
\item How to {\bf produce} DDH files from ARPEGE AROME HARMONIE models ? 
\item Flexible DDH.
\item How to {\bf use} DDH files for budget and plots ? The ddhtoolbox.
\item Where to find DDH documentation ? Who to contact ?
\item Conclusions
%\item <3-> 
\end{itemize}
\end{block}
\end{columns}
\end{textblock*}

\logos

\end{frame}
%---------------------------------
% ajouter sur la slide suivante :
%---------------------------------
\addtocounter{framenumber}{-1}

%============================================================================
\subsection{Why use DDH ?}
\begin{frame}
\frametitle{\vspace{-0.3cm} \centerline{\LARGE Why use DDH ?} }

\vspace{-1.0cm}
\begin{exampleblock}{\normalsize{Objectives :}}
\begin{itemize}
 \item Establish balance equations for the model's prognostic variables over a defined domain, 
to improve understanding of the model's dynamics and physical interactions. 
 $\Rightarrow$ Contribute to the refinement of physical parameterizations.
 \item Provide the value of diagnostic fields (cloud cover, pressure, etc.).
\end{itemize}
\end{exampleblock}


\begin{block}{\normalsize{Output:}}
\begin{itemize}
 \item \textbf{Budgets} : mass (PP), water vapour (QV), kinetic energy (KK), angular momentum (A1,A2,A3), entropy (SS), potential energy (EP), temperature (CT), etc.
 \item \textbf{Domains} : the whole model domain, zonal bands, limited area domains (rectangles, polygons), isolated grid points. 
\end{itemize}
\end{block}
\end{frame}

%============================================================================

\subsection{Budget equation}
\begin{frame}
\frametitle{\normalsize Generic form of budget equations}
\vspace{-0.25cm}
\begin{block}{\small Budget equation in $\eta$ coordinate for the variable $X$ :}
%  
\small
\begin{align*}
\frac{1}{g}\frac{\dr}{\dr t}\lp \frac{\dr \pi}{\dr\eta}X \rp & = 
\frac{1}{g}\frac{\dr \pi}{\dr\eta}\sum_{\mathpzc{D}} \mathcal{T}_{\mathpzc{D}} 
- \sum_{\varphi} \frac{\dr \mathcal{F}_{\varphi}}{\dr\eta}
+ \mathcal{E}
\end{align*}
%
\begin{itemize}
\item $\mathcal{T}_{\mathpzc{D}}$ : Dynamical tendencies, e.g horizontal or vertical advection, pressure gradient,  $\ldots$
%
\item $\mathcal{F}_{\varphi}$ : Physical fluxes, e.g [radiation, convection, turbulence, precipitation, $\ldots$], from parametrizations.
%
\item $\mathcal{E}$ : Residual, unexplained by physical or dynamical explicit terms, e.g [horizontal diffusion, semi-implicit correction, ...]
\end{itemize}
%
\end{block}
\end{frame}
%============================================================================
%============================================================================

\subsection{Thermal energy equation}
\begin{frame}
\frametitle{{\normalsize An example :} }

\begin{block}{Thermal energy budget $c_p T$}
\footnotesize
\begin{align}
  \DerP{}{t}\lp\re c_p T\rp &=
  - \divi{\eta}\lp\re c_p T \vec{v}\rp
  - \DerP{}{\eta}\lp\re c_p T \dot{\eta}\rp
  + \re R T \frac{\omega}{p}\nonumber \\
  &  +\DerP{}{\eta}\lb F_{c_p T} + \fcptp{}
  + \fpl{} T \lc c_l - \cp{a} (1-\delta_m)\rc 
  + \fpn{} T \lc c_n - \cp{a} (1-\delta_m)\rc\rb\nonumber\\
  & 
  +\delta_m\fp{}\DerP{(\Phi+\frac{u^2+v^2}{2})}{\eta}
  -\vec v\cdot\DerP{\vec F_v^{\varphi}}{\eta}\label{Bilh}
\end{align}
\vspace*{0.5cm}\\
avec $r_{\eta}=-(1/g)\dr_{\eta}\pi$
\end{block}


\logos

\end{frame}

%============================================================================
\subsection{ARPEGE global mean budget}
\begin{frame}
\frametitle{\vspace{0.2cm} \hspace{3.5cm} {\LARGE Budget equations} }

\begin{textblock*}{11.cm}(0.5cm,1.9cm)
\includegraphics<1->[width=11.cm]{zddhb.CT.OPER.graph.doc.eps}
\end{textblock*}

\begin{textblock*}{12.5cm}(1mm,1.3cm)
\begin{columns}
\column{11.5cm}
\begin{block}<-1>{}
\begin{center}
ARPEGE global mean temperature budget
\end{center}
\end{block}
\end{columns}
\end{textblock*}

\logos

\end{frame}

%============================================================================
\subsection{Budget equations}
\begin{frame}
\frametitle{\vspace{0.2cm} \hspace{3.5cm} {\LARGE Budget equations} }

\begin{textblock*}{12.cm}(0.5cm,2.7cm)
\includegraphics<1->[width=12.cm]{yann.eps}
\end{textblock*}

\begin{textblock*}{12.5cm}(1mm,1.3cm)
\begin{columns}
\column{11.5cm}
\begin{block}<-1>{}
\begin{center}
AROME temperature budget over the Chamonix valley. \\ Source Yann Seity.
\end{center}
\end{block}
\end{columns}
\end{textblock*}

\logos

\end{frame}

%============================================================================
\subsection{Meridional or zonal circulations, in cross-sections}
\begin{frame}
\frametitle{\vspace{-0.3cm} \centerline{\Large Meridional or zonal circulations, in cross-sections} }

\begin{textblock*}{2.7182818449cm}(-0.8cm,1.8cm)
\includegraphics<1->[width=18.cm]{DHFZOFCST+0072_xyuvc_1_doc_svg_HR.eps}
\end{textblock*}


\end{frame}

%%%%%%%%%%%%%%%%%%%%%%%%%%%%%%%%%%%%%%%%%%%%%%%%%%%%%%%%%%%%%%%%%%%%%%%%%%%%%
\section{Producing DDH files}


%============================================================================
\subsection{Producing DDH files}
\begin{frame}
\frametitle{\vspace{-0.3cm} \centerline{\LARGE Producing DDH files } }

\begin{textblock*}{2.7182818449cm}(0.8cm,1.8cm)
\includegraphics<1->[width=9.cm]{production.eps}
\end{textblock*}

\begin{textblock*}{2.7182818449cm}(11cm,4.cm)
\includegraphics<1->[width=5.cm]{horloge-indus-a-rouages-effet-rouille-d125-500-13-29-138240_1.eps}
\end{textblock*}

\begin{textblock*}{2.7182818449cm}(10.6cm,1.8cm)
\includegraphics<1->[width=5.5cm]{BullSequanaX-key-visual-538-303.eps}
\end{textblock*}

\begin{textblock*}{2.7182818449cm}(1.cm,6cm)
\includegraphics<1->[width=1.1cm]{facile.eps}
\end{textblock*}


\end{frame}


%============================================================================
\subsection{History}
\begin{frame}
\frametitle{\vspace{-0.3cm} \centerline{\LARGE Some history } }
\vspace{-.7cm}
\begin{block}{\normalsize{Developing DDH inside our models}}
{\small
\begin{description}
  \item [1991:] Initial analysis and coding of the {\tt DDH} software (Alain Joly).
  \item [1992:] Introduction of entropy and kinetic energy budgets (Martin Janousek, Jean-Marcel Piriou).
  \item [1993:] Introduction of relative humidity, liquid and ice water (Jean-Marcel Piriou).
  \item [1997:] New surface fields, IFS model (Pedro Viterbo).
  \item [2000:] Diagnose surface "tiles", IFS model (Christian Jakob).
  \item [2006:] Extract AROME physical data flow, interface to DDH routines (Tomislav Kovacic).
  \item [2006:] Write AROME DDH documentation (Tomislav Kovacic).
  \item [2007-11:] Rewrite the budget tool "ddhb",
    still based on the "ddhi" and "ddht" existing ones (Alex Deckmyn, Jean-Marcel Piriou, Tomas Kral).
\end{description}
}
\end{block}
\end{frame}

%============================================================================
\subsection{History}
\begin{frame}
\frametitle{\vspace{-0.3cm} \centerline{\LARGE Some history } }
\vspace{-.7cm}
\begin{block}{\normalsize{Developing DDH inside our models}}
{\small
\begin{description}
  \item [2008-07:] Create the ddhtoolbox, write its documentation (Jean-Marcel Piriou).
  \item [2009:] Interface AROME physics with DDH, new flexible dataflow
    for AROME, ALADIN and ARPEGE (Olivier Rivière).
  \item [2018:] Dynamical DDH: semi-lagrangian advection, horizontal diffusion, semi-implicit (Fabrice Voitus).
  \item [2019:] Flexible DDH: adding a new field in calling NEW\_ADD\_FIELD\_3D (Fabrice Voitus).
  \item [2023:] A new section was written in the ddhtoolbox documentation, to learn step by step, from examples, the use of the DDH tools (ddhi, ddhr, ddht, etc) (Jean-Marcel Piriou, Panu Maalampi).
  \item [2023:] The dd2gr plot tool is now part of the ddhtoolbox, to allow an easier start with DDH use (Jean-Marcel Piriou).
\end{description}
}
\end{block}
\end{frame}
%
%============================================================================

\subsection{Récupération des tendances}

\begin{frame}[fragile]
\frametitle{\vspace{-0.3cm} \centerline{\LARGE Get tendencies, fluxes into DDH variables } }
\vspace{-0.45cm}
\begin{alertblock}{} 
\footnotesize
\begin{itemize}
\item A call to NEW\_ADD\_FIELD\_3D or NEW\_ADD\_FIELD\_2D via the module \textbf{ddh\_mix} anywhere under the \textbf{cpg} or \textbf{cpglag} routines allows the storage of a variable, a tendency, or a flux as a vertical profile inside DDH files.
\end{itemize}
\end{alertblock}

\begin{exampleblock}{} 
\begin{footnotesize}
\begin{itemize}
\item Store the X variable:
\begin{verbatim}
    ZTMPVAR(KIDIA:KFDIA,1:KFLEV)=&
      & PDELP(KIDIA:KFDIA,1:KFLEV)*PQ(KIDIA:KFDIA,1:KFLEV)
    CALL NEW_ADD_FIELD_3D(YDMODEL%YRML_DIAG%YRMDDH,ZTMPVAR,'VQV',YDDDH)
\end{verbatim}
\item Store the Tendency of X from a given process:
\begin{verbatim}
     ZTNDKKT(:,:) = ...
     CALL NEW_ADD_FIELD_3D(YDMODEL%YRML_DIAG%YRMDDH,ZTNDKKT,'TKKDISSTUR',YDDDH)
\end{verbatim}
 \item Store the Flux of X from a given process:
\begin{verbatim}
     CALL NEW_ADD_FIELD_3D(YDMODEL%YRML_DIAG%YRMDDH,PDIFTQ,'FQVTUR',YDDDH)
\end{verbatim}
\end{itemize}
\end{footnotesize}

\end{exampleblock}

\end{frame}
%============================================================================


%============================================================================
\subsection{Capuccino}

\begin{frame}[fragile]
\frametitle{{\normalsize An example : CAPUCCINO} }
\vspace*{-0.25cm}
\begin{exampleblock}{}
\begin{tiny}
 \begin{verbatim} 
  SUBROUTINE DIAG_CAPUCCINO(KIDIA,KFDIA,KLEV,PINGREDIENTS,PDOSAGE,YDLDDH,YDMDDH,YDDDH)
  
  USE DDH_MIX   , ONLY : ADD_FIELD_3D, NEW_ADD_FIELD_3D
  USE YOMMDDH, ONLY  : TMDDH
  USE YOMLDDH, ONLY  : TLDDH     
  ....
  TYPE(TLDDH)       ,INTENT(IN)    :: YDLDDH ! ddh useful logicals
  TYPE(TMDDH)       ,INTENT(INOUT) :: YDMDDH ! ddh useful metadata 
  TYPE(TYP_DDH)     ,INTENT(INOUT) :: YDDDH  ! ddh superstructure  
  ....
  DO JLEV=1,KLEV
     DO JLON = KIDIA,KFDIA      
        ZCOFFEE(JLON,JLEV)      = .... 
        ZMILK(JLON,JLEV)        = ....
        ZCHOCOLATE(JLON,JLEV)   = ....
     ENDDO
  ENDDO   
  ....
  IF (YDLDDH%LFLEXDIA).AND.(YDLDDH%LDDH_OMP) THEN
     CALL NEW_ADD_FIELD_3D(YDMDDH,ZCOFFEE,'FCO_COFFEE',YDDDH)
     CALL NEW_ADD_FIELD_3D(YDMDDH,ZMILK,'CO_MILK',YDDDH,CDTYPE='F')
     CALL NEW_ADD_FIELD_3D(YDMDDH,ZCHOCOLATE,'FCO_CHOCO',YDDDH)
  ENDIF
  ....
  END  SUBROUTINE DIAG_CAPUCCINO
  \end{verbatim}  
  \end{tiny}
\end{exampleblock}
\end{frame}

%============================================================================
\subsection{User Namelist}

\begin{frame}[fragile]
\frametitle{{\normalsize User Namelist : an example} }
\vspace*{-0.4cm}
\begin{exampleblock}{}
\begin{footnotesize}
\begin{verbatim} 
 &NAMDDH                                                                                            
  BDEDDH(1,1)=4., ! domain : a single grid point.
  BDEDDH(2,1)=1., ! virtual plane.
  BDEDDH(3,1)=1.374, ! longitude.
  BDEDDH(4,1)=43.575, ! latitude.

  BDEDDH(1,2)=3., ! domain : rectangle in lat-lon.
  BDEDDH(2,1)=1., ! virtual plane.
  BDEDDH(3,2)=0.0, ! corner 1: longitude.
  BDEDDH(4,2)=45.0, ! corner 1: latitude.
  BDEDDH(5,2)=0.1, ! corner 2: longitude.
  BDEDDH(6,2)=44.8, ! corner 2: latitude.

  LHDDOP=.TRUE., ! activates limited aread domains computation.
  LHDEFD=.TRUE., ! activates limited aread domains output on files.
  LHDHKS=.TRUE., ! activates the standard DDH fields package.
  LFLEXDIA=.TRUE., ! activates flexible DDH.
  LDDH_OMP=.TRUE., ! activates thread-safe DDH.
 /
\end{verbatim}  
\end{footnotesize}
\end{exampleblock}
\end{frame}
%%%%%%%%%%%%%%%%%%%%%%%%%%%%%%%%%%%%%%%%%%%%%%%%%%%%%%%%%%%%%%%%%%%%%%%%%%%%%
\section{Using DDH files}

%============================================================================
\subsection{Planche d'annonce}
\begin{frame}
%\frametitle{\vspace{0.2cm} \hspace{3.5cm} {\LARGE Circulations méridiennes} }
\frametitle{\vspace{-0.3cm} \centerline{\LARGE Using DDH files} }

\begin{textblock*}{2.7182818449cm}(0.8cm,1.8cm)
\includegraphics<1->[width=12.cm]{enfants.eps}
\end{textblock*}

\logos

\end{frame}
%============================================================================
\subsection{ddhtoolbox}
%============================================================================
\begin{frame}
%\frametitle{\vspace{0.5cm} \hspace{3.5cm} {\LARGE Sommaire} }
\frametitle{\vspace{0.3cm} \centerline{{\LARGE ddhtoolbox}} }

\begin{textblock*}{11.cm}(1cm,1.5cm)
\includegraphics<1->[width=11.cm]{synoptic.eps}
\end{textblock*}

\logos

\end{frame}


%============================================================================
\subsection{Emagramme}
%============================================================================
\begin{frame}
%\frametitle{\vspace{0.5cm} \hspace{3.5cm} {\LARGE Sommaire} }
\frametitle{\vspace{-0.3cm} \centerline{{\LARGE Plotting Tphigrams}} }

\begin{textblock*}{2.7182818449cm}(0.3cm,2.5cm)
\includegraphics<1->[width=8.cm]{theta_toulouse.eps}
\end{textblock*}

\begin{textblock*}{2.7182818449cm}(8.3cm,2.5cm)
\includegraphics<1->[width=8.cm]{theta_tropiques.eps}
\end{textblock*}


\begin{textblock*}{13.5cm}(1.5cm,1.5cm)
\begin{columns}
\column{13.5cm}
\begin{block}<2->{}
% \Huge \huge \ LARGE \Large \large
% \normalsize \small \footnotesize \scriptsize \tiny
{\footnotesize
\begin{itemize}
\item get hendrix:/home/m/mxpt/mxpt001/vortex/arpege/4dvarfr/OPER/2019/09/17 \\
/T0000P/forecast/ddh.arpege-zonal.tl1798-c22.ddhpack.tgz 
\item tar tvf
\item ddhr -ddd DHFZOFCST+0072 \textbf{\green{get latitude, longitude of each domain}}
\item ddht -cEXTRAIT\_DOMAIN -E13-18 -1DHFZOFCST+0018 -sDHFZOFCST+0018.Tropiques.lfa \textbf{\green{extract some domains}}
\item ddhmh DHFZOFCST+0018.Tropiques.lfa \textbf{\green{make horizontal mean}}
\item ddh2scm DHFZOFCST+0018.Tropiques.lfa.mh \textbf{\green{convert into a single-column-type MUSC file}}
\item ms DHFZOFCST+0018.Tropiques.lfa.mh.Dom001.Var\_fin.scm \textbf{\green{model sounding}}
\item dd2gr DHFZOFCST+0018.Tropiques.lfa.mh.Dom001.Var\_fin.scm.ms.tmp.theta.doc tphi.svg \textbf{\green{plot the sounding into a SVG file}}
\item convert tphi.svg tphi.png \textbf{\green{convert SVG file into a PNG (or EPS, PDF) file}}
%\item <3-> 
\end{itemize}
}
\end{block}
\end{columns}
\end{textblock*}



\end{frame}

%============================================================================
\subsection{Hovmöller diagram time-height-value}
\begin{frame}
\frametitle{\vspace{-0.3cm} \centerline{{\LARGE Hovmöller diagram time-height-value}} }

\begin{textblock*}{2.7182818449cm}(2.cm,1.7cm)
\includegraphics<1->[width=11.5cm]{VNT1.doc.eps}
\end{textblock*}
\end{frame}

%============================================================================
\subsection{Budget, 21 November in Madrid, 08 UTC}
\begin{frame}
\frametitle{\vspace{-0.3cm} \centerline{{\LARGE Budget, 21 November in Madrid, 08 UTC}} }

\begin{textblock*}{2.7182818449cm}(0.cm,2.3cm)
\includegraphics<1->[width=16.cm]{DHFDLFCST+0032.instantane.comb.eps}
\end{textblock*}
\end{frame}

%============================================================================
\subsection{Budget, 21 November in Madrid, 09 UTC}
\begin{frame}
\frametitle{\vspace{-0.3cm} \centerline{{\LARGE Budget, 21 November in Madrid, 09 UTC}} }

\begin{textblock*}{2.7182818449cm}(0.cm,2.3cm)
\includegraphics<1->[width=16.cm]{DHFDLFCST+0033.instantane.comb.eps}
\end{textblock*}
\end{frame}

%============================================================================
\subsection{Budget, 21 November in Madrid, 10 UTC}
\begin{frame}
\frametitle{\vspace{-0.3cm} \centerline{{\LARGE Budget, 21 November in Madrid, 10 UTC}} }

\begin{textblock*}{2.7182818449cm}(0.cm,2.3cm)
\includegraphics<1->[width=16.cm]{DHFDLFCST+0034.instantane.comb.eps}
\end{textblock*}
\end{frame}

%============================================================================
\subsection{Budget, 21 November in Madrid, 11 UTC}
\begin{frame}
\frametitle{\vspace{-0.3cm} \centerline{{\LARGE Budget, 21 November in Madrid, 11 UTC}} }

\begin{textblock*}{2.7182818449cm}(0.cm,2.3cm)
\includegraphics<1->[width=16.cm]{DHFDLFCST+0035.instantane.comb.eps}
\end{textblock*}
\end{frame}

%============================================================================
\subsection{Budget, 21 November in Madrid, 12 UTC}
\begin{frame}
\frametitle{\vspace{-0.3cm} \centerline{{\LARGE Budget, 21 November in Madrid, 12 UTC}} }

\begin{textblock*}{2.7182818449cm}(0.cm,2.3cm)
\includegraphics<1->[width=16.cm]{DHFDLFCST+0036.instantane.comb.eps}
\end{textblock*}
\end{frame}
%%%%%%%%%%%%%%%%%%%%%%%%%%%%%%%%%%%%%%%%%%%%%%%%%%%%%%%%%%%%%%%%%%%%%%%%%%%%%
\section{DDH Documentation }

%============================================================================
\subsection{DDH Documentation}
\begin{frame}
\frametitle{\vspace{-0.3cm} \centerline{\LARGE DDH Documentation} }

\vspace{-1.0cm}
\begin{exampleblock}{} 
\begin{itemize}
\item {\bf GitHub of the ddhtoolbox}: \href{https://github.com/UMR-CNRM/ddhtoolbox}{https://github.com/UMR-CNRM/ddhtoolbox} 
\item {\bf DDH documentation} : on the above GitHub : documentation/ddh.pdf : all budget equations, user namelist, DDH file structure, DDH tools (ddhtoolbox), examples of use
\end{itemize}
\end{exampleblock}

\begin{block}<2->{\normalsize{Hotline exceptional level :) }}
\begin{itemize}
 \item {\bf Fabrice Voitus} for producing DDH files from 3D or 1D models ARPEGE, AROME.
 \item {\bf Jean-Marcel Piriou, Marie Mazoyer} for the ddhtoolbox use on PC. 
 \item {\bf Panu Maalampi, Marvin Kähnert, Daniel Martin, Emily Gleeson} for producing DDH files from 3D or 1D models HARMONIE, and for the ddhtoolbox use on PC. 
 \item {\bf Enjoy DDH use} and create a wonderful model physics !
\end{itemize}
\end{block}
\end{frame}


%============================================================================
\end{document}
